%%%%%
%
\section{Sommation de plusieurs "colonnes" d'un param\`etre 1D}

\itemstyle{Op\'erateur~:} {\sf Som\_Col}
\index{Som\_Col}

%%%%%OBJET%%%%
\itemstyle{Objet~:}

Sommation de plusieurs colonnes cons\'ecutives d'un
param\`etre 1D. 

%%%%%ARGUMENTS%%%%
\bigskip
%%%%%ENTREES%%%%
\begin{tabarg}
\multicolumn{3}{l}{\titretab {Entr\'ees}}\\
\hline
$X$ & Vecteur, \texttt{Int16[$D$]} & Param\`etre 1D d'entr\'ee (Unit\'e quelconque)\\
$D$ & Coeff., \texttt{Int} & Nombre d'\'el\'ements d'un \'echantillon de $X$ \\
$C_{deb}$ & Coeff., \texttt{Int} & Num\'ero de la premi\`ere colonne (de $0$ � $D-1$)\\
$C_{fin}$ & Coeff., \texttt{Int} & Num\'ero de la derni\`ere colonne (de $0$ � $D-1$ et $\ge C_{deb}$ )\\
\end{tabarg}

\medskip
%%%%%SORTIES%%%%
\begin{tabarg}
\multicolumn{3}{l}{\titretab {Sorties}}\\
\hline
$Z$ & Vecteur, \texttt{float} & Somme (m\^eme unit\'e que $X$)\\
\end{tabarg}


%%%%%FORMULE%%%%
\itemstyle{Formule}

%%som_15vect
\[
Z = \sum_{k=C_{deb}}^{C_{fin}} X_k
\]


%%%%COMENTAIRE
\itemstyle{Commentaire~:}

La premi\`ere colonne porte le num\'ero $0$.\\
Pour le premier argument, on peut utiliser la syntaxe : \texttt{\$Arg1.array\_dim}.

%%%%%AUTEUR%%%%
\itemstyle{Auteur~:}
CNRM/GMEI/TRAMM

%%%%%REFERENCE%%%%
\itemstyle{R\'ef\'erences~:}
